\clearpage\thispagestyle{empty}
\pdfbookmark[0]{Quellen zur Rohde-Rezension von 1805}{rohde1805-sources}
\begin{center}\textbf{Quellen zur Rohde-Rezension von 1805}\end{center}
\noindent Diese KI-Übertragung ins heutige Deutsch folgt dem Abteilungstitel und der Rezension auf S.~483--485 der \href{https://gdz.sub.uni-goettingen.de/id/PPN236007467}{\textit{Werke}, Band~VI} (Göttingen, 1874). Carl Friedrich Gauß bespricht die Erwiderung des Hauptmanns Rohde auf eine Rezension seiner Abhandlung über die Massen der Planeten. Der folgende Briefauszug von Lindenau nennt Gauß als den Rezensenten.
\par\smallskip\noindent Die Rezension erschien als Teil V von XXXIX, \textit{Critik und Anticritik}, in Franz Xaver von Zachs \textit{Monatlicher Correspondenz}, Band~XII, Oktober 1805, S.~\href{https://api.digitale-sammlungen.de/iiif/image/v2/bsb10538604_00429/full/1500,/0/default.jpg}{413}--\href{https://api.digitale-sammlungen.de/iiif/image/v2/bsb10538604_00435/full/1500,/0/default.jpg}{419}. Die Werkausgabe setzt anstelle der Punkte im Erstdruck Striche.
\par\smallskip\noindent Die beiden französischen Zitate in Bemerkung I bedeuten \emph{„um die Ernte von allen Seiten zu beschleunigen, für die die Kräfte eines einzigen Mannes nicht ausreichen, meine aber überhaupt nicht“} und \emph{„angesichts der großen Zahl neuer analytischer Kombinationen usw.“} Die Bemerkung in Klammern dazwischen stammt von Gauß. KI-Übertragung und KI-Satz.
\clearpage
